\subsection{Growth, the error budget and the first section}

Fix all compact classes and the finite cones in
Theorem~\ref{thm:detuned-output-admission}.  Reduce the graph outer radius first, and
then take its compatible reset cutoff, so that
\[
 0<\delta_{\rm reset}<\delta_{\rm g}<\delta_{\rm call},\qquad
 q_{\rm g}:=(1+L_\Gamma)q<1,
\]
where $q<1$ is the full logarithmic contraction in
\eqref{gr:contraction} and $L_\Gamma$ is the logarithmic Lipschitz
constant of the graph on the outer interval.  Choose a fixed
$\Delta<\delta_{\rm reset}$ below every inner input cutoff.  It is
also chosen so small that
\begin{equation}
 c_D\Delta^{-1/4}>2,
 \qquad c_D\delta^{3/4}\le\delta^+\le C_D\delta^{3/4}
 \quad(0<\delta\le\Delta),
 \label{path:growth-bounds}
\end{equation}
and the full-prefix tubes below lie in a common neighborhood
$\mathcal U_G$ with $p_\phi\omega<-a/2$.  These choices precede the
choice of the observation section.

Write $\mathcal B_\lambda(d)$ for the value of the dependent
coordinate $\eta_B=(-b-c_\lambda d)/(c_\lambda d^{4/3})$
at the graph state $(d,\Gamma_\lambda(d))$. It is obtained by
substituting that state into the physical inverse \eqref{loc:inverse}
with $\eps=r=0$ and evaluating $b$. In the notation of
Appendix~\ref{sec:detuned-return}, this is
$\mathcal B_{0,0}(d,\Gamma_\lambda(d))$; the corresponding
Lipschitz estimate is \eqref{ret:dependent-Lipschitz}.

\begin{lemma}[Stopped cascade and its complete state budget]
\label{path:stopped-cascade}
For each fixed $0<h\le\Delta$, the actual returns from $\TD_3$ can
be generated, on the original trajectory, up to
\[
 N_h=\min\{n\ge3:\delta_n\ge h\}.
\]
This index is finite, tends to infinity as $\eps\downarrow0$, and
every input used before it is below the reset and inner cutoffs.
The terminal output remains in the graph outer domain.  The complete
cone, inverse branch, negative memory and strict first-event margins
are retained.  Put
\[
 e_n=\eps\delta_n^{-3/2},\quad
 \omega_n=e_n+r/\delta_n,\quad
 B_n=\|\xi_n^r-\Gamma_\lambda(\delta_n)\|_\infty.
\]
For $3\le m\le N_h$,
\begin{align}
 B_m&\le q_{\rm g}^{m-3}B_3+
       C\sum_{n=3}^{m-1}q_{\rm g}^{m-1-n}\omega_n,
       \label{path:graph-budget}\\
 \left|\log\frac{\delta_{n+1}}{F_\lambda(\delta_n)}\right|
       &\le qB_n+C\omega_n\quad(3\le n<N_h),
       \label{path:same-base-output}\\
 \Omega_\eps:=\sup_{3\le n<N_h}\omega_n
       &\le C(\eps^{37/64}+\eps^{7/32})\longrightarrow0.
       \label{path:uniform-defect}
\end{align}
If an actual and a limiting chain are compared in full logarithmic
state, their distance satisfies, through the first output at which
either input chain leaves the common inner domain,
\begin{equation}
 E_m\le q^{m-3}E_3+
          C\sum_{n=3}^{m-1}q^{m-1-n}\omega_n.
 \label{path:two-chain-budget}
\end{equation}
At any sequence of positive-scale inputs whose indices tend to
infinity, the dependent acceleration coordinate also tends to its
graph value.  Thus this convergence concerns the complete normalized
touchdown state, including memory.
\end{lemma}

\begin{proof}
Proposition~\ref{loc:entrance} gives the actual entrance
\begin{equation}
 \delta_3\asymp\eps^{9/32},\qquad
 \eps\le C_0\delta_3^2,\qquad
 r\le A\delta_3^{4/3},\qquad
 x_3=\tfrac9{32}\log(1/\eps)+O(1),
 \quad x=\log(K_\lambda^4/\delta).
 \label{path:entrance}
\end{equation}
The admissibility of the returned state in
Theorem~\ref{thm:detuned-output-admission} and
\eqref{path:growth-bounds} give strict scale growth.  Since $\eps$ and $r=k\sqrt\eps$ stay
fixed on each trajectory, both admissibility inequalities propagate
to the next input.  The return estimate is applied again when its new input
is below $h\le\Delta$.  All events are the
actual first release and subsequent first touchdown, with positive
gap between them.  If there were infinitely many inputs below $h$,
\eqref{path:growth-bounds} would give $\delta_n\ge2^{n-3}\delta_3$,
a contradiction.  Conversely every fixed finite index has
$\delta_n\to0$, by \eqref{path:entrance} and the growth upper bound.
Consequently $N_h\to\infty$ uniformly on the compact classes.

The direct value comparison in
Theorem~\ref{thm:detuned-value-bridge} is made at identical values of the
independent normalized touchdown coordinates and gives a logarithmic output defect at most
$C\omega_n$.  Integrating the derivative bound for the limiting
map along the logarithmic joining segment gives
$E_{n+1}\le qE_n+C\omega_n$, and hence
\eqref{path:two-chain-budget}.  Only the limiting map is
differentiated.  To compare with the graph, take its input at the
actual base $\delta_n$.  Graph invariance and the same argument give
\[
 \max\left\{
 \left|\log\frac{\delta_{n+1}}{F_\lambda(\delta_n)}\right|,
 \|\xi_{n+1}^r-\Gamma_\lambda(F_\lambda(\delta_n))\|_\infty
 \right\}\le qB_n+C\omega_n.
\]
Both output bases are in the outer graph interval.  Correcting the
graph base there costs $L_\Gamma(qB_n+C\omega_n)$, so
$B_{n+1}\le q_{\rm g} B_n+C\omega_n$.  Iteration proves
\eqref{path:graph-budget}.  The bounded entrance error and the two
geometric sums give bounds independent of the number of returns.
Scale growth also gives
\[
 \Omega_\eps\le\eps\delta_3^{-3/2}+r/\delta_3
           \le C(\eps^{37/64}+\eps^{7/32}),
\]
which proves \eqref{path:uniform-defect}.

To compare complete states, recover the acceleration row from the
same physical inverse.
At a common base $d$, the actual and graph arguments differ by
\begin{equation}
 \begin{split}
 \Delta p&=0,\qquad
 \Delta u_0=d_\lambda d^{4/3}
       (\eta_U^r-\Gamma_U(d))+r,\\
 \Delta v&=-r_0d^{7/3}(\eta_W-\Gamma_W(d)),\qquad
 \Delta y=\eps d^2\zeta.
 \end{split}
 \label{path:inverse-arguments}
\end{equation}
Smoothness of that inverse and division by $c_\lambda d^{4/3}$ give
\begin{equation}
 |\eta_B-\mathcal B_\lambda(d)|
 \le C\{B_n+r/d^{4/3}+\eps d^{2/3}|\zeta|\}.
 \label{path:dependent-row}
\end{equation}
On any fixed positive-scale compact interval this tends to zero
when $n\to\infty$ and $\eps\to0$.  The memory component of the graph vanishes,
while the actual values at touchdown remain exactly
$y=z=\eps d^2\zeta$.  No normal coordinate is reset.
\end{proof}

We record the prefix estimates needed to turn this state budget
into a first-crossing and path argument.  The contact estimates
\eqref{con:prefix}, followed by Lemma~\ref{lem:paired-free-normalization}, give,
for every admitted input $\delta\le\Delta$ and every actual prefix,
\begin{equation}
 \begin{gathered}
 \sup_{\rm contact}|X-G|\le C\delta,
 \qquad t_L-t_{\rm TD}\le C\eps/\sqrt\delta,\\
 \sup_{\rm contact}(|y|+|z|)\le C\eps\delta^{3/2}.
 \end{gathered}
 \label{path:contact-tube}
\end{equation}
\begin{equation}
 \begin{gathered}
 T:=t_+-t_L\le H\delta^{3/4},\qquad
 |X(t_L+\ell)-G|\le C(\delta+\ell),\\
 z(t_L+\ell)=z_Le^{-\ell/\eps}<0
 \qquad(0\le\ell\le T).
 \end{gathered}
 \label{path:free-tube}
\end{equation}
In particular, since $|b_L|\le C\delta$ and
$b(t_L+\ell)=O(\delta+\ell)$, integration of $\dot v=b$ and
$\dot y=v$ gives
\begin{equation}
 \begin{split}
 |z|/\eps&\le C\delta^{3/2},\qquad
 |v|\le C(\delta^{3/2}+\delta\ell+\ell^2),\\
 |y|&\le C(\eps\delta^{3/2}+\delta^{3/2}\ell
                       +\delta\ell^2+\ell^3)\le C\delta^{9/4}.
 \end{split}
 \label{path:free-normal-tube}
\end{equation}
The contact velocity estimate uses $e/\sqrt\delta\le C_0$ in
the weighted compression bounds.  Thus these estimates remain valid
without a positive lower bound on $e$.

\begin{lemma}[Original preparation and the first physical section]
\label{path:first-section}
The whole actual curve from time zero through $\TD_3$ tends to the
prescribed incoming arc $I_{\lambda,\nu}$ in the completed ordered
distance.

There is a common $s_0>0$ such that for fixed $0<s\le s_0$ the
first crossing $t_s$ of $p=-as$ from the original initial state
exists and is transverse.  With $N=N_s$,
\begin{equation}
 t_{N-1}<t_s\le t_N,
 \qquad c_s:=(s/C_D)^{4/3}\le d_\eps=\delta_{N-1}<s,
 \label{path:section-location}
\end{equation}
and the complete same-base defect
\begin{equation}
 \kappa_\eps=
 |\xi_{N-1}^r-\Gamma_\lambda(d_\eps)|+
 |\eta_{B,N-1}-\mathcal B_\lambda(d_\eps)|\longrightarrow0.
 \label{path:last-input}
\end{equation}
All statements here are uniform on the fixed compact label classes.
\end{lemma}

\begin{proof}
Fix a small incoming cutoff $\eta>0$. Proposition~\ref{loc:entrance}
gives convergence to the corresponding rigid prefix. Transverse
alignment on that closed subarc gives an increasing pairing.
The remaining rigid tail lies in $G+B_{C\eta}$, while the proposition
puts the remaining original prefix through $\TD_3$ in
$G+B_{C\eta+o(1)}$. Pair these tails in their order and then let
$\eta\downarrow0$; this proves the ordered incoming limit.

For the cascade up to $N_s$, take the supremum of the local bounds:
\begin{equation}
 \sup_{[t_3,t_N]}|X-G|\le Cs^{3/4},\qquad
 \sup_{[t_3,t_N]}|z|/\eps\le Cs^{3/2}.
 \label{path:common-tube}
\end{equation}
There is no sum over returns.  Choose $0<s_0<\Delta/2$ so this tube lies in
$\mathcal U_G$.  The exact identity
\begin{equation}
 \dot p=p_\phi(\phi)\omega<-a/2
 \label{path:strict-p}
\end{equation}
then holds on every contact and free prefix and at their joins.
The entrance has $\delta_3<s$, whereas $\delta_N\ge s$, giving a
unique crossing between the stated events.  On the original closed
incoming prefix $p$ is positive; the part from its small cutoff to
the threshold also has positive $p$.  From the threshold through
$\TD_3$, Proposition~\ref{loc:entrance} gives $\sup|p|=o(1)$.  For fixed
$s$, this excludes an earlier crossing of $-as$.  Hence the crossing
is the first one from the original initial data.  Only after the
tube has been established, integration of \eqref{path:strict-p}
also gives $t_s-t_3\le2(s-\delta_3)$.

The last growth upper bound gives $d_\eps\ge c_s$.  The last-input
index tends to infinity, so \eqref{path:graph-budget} and
\eqref{path:dependent-row} give \eqref{path:last-input} on
$[c_s,s]$.
\end{proof}

\begin{corollary}[Localization of an actual contact stop]
\label{cor:contact-location}
Fix $0<s\le s_0$. If an actual touchdown generated by the return estimates has
$\delta_j\in[s/2,2s]$ and the first section time satisfies
$t_j<t_s\le t_{L,j}$, then
\begin{equation}
 0<s-\delta_j\le C_s\eps,\qquad
 0<t_s-t_j\le C_s\eps.
 \label{eq:contact-stop-location}
\end{equation}
The constant is uniform on the prescribed compact label classes.
\end{corollary}

\begin{proof}[Proof of Corollary~\ref{cor:contact-location}]
On the fixed band $[s/2,2s]$, \eqref{path:contact-tube} gives
$t_{L,j}-t_j\le C_s\eps$. Since $\dot p$ is bounded and satisfies
\eqref{path:strict-p} throughout this contact,
\[
 0<s-\delta_j=-\frac1a\int_{t_j}^{t_s}\dot p\,\dd t
       \le C_s(t_s-t_j)\le C_s\eps.
\]
The same contact bound gives the asserted time estimate.
\end{proof}


For a first section time during contact $j$, Lemma~\ref{path:first-section}
gives $j=N_s-1$ and $\delta_j\in[c_s,s)$. The same contact-time
argument on this fixed positive band gives $s-\delta_j=O_s(\eps)$,
so the input eventually lies in $[s/2,2s]$.
Thus \eqref{eq:contact-stop-location} applies to every actual contact stop,
without assuming the band in advance.

\subsection{Sharp comparison from the actual third touchdown}

Fix $\lambda$ and a compact auxiliary class $\mathfrak K$ within
the standing ranges. Put $A_0=9/32$, $\rho_c=3/4$ and
$L=\log(1/\eps)$. The scalar phase contracts by $\rho_c$;
the complete-state contraction is $q=7/8$.
The entrance state is already $O(\eps^{3/32})$ from the graph at
its own base, so its complete-state budget resolves the residual
entrance constant through all required returns.

\begin{lemma}[Same-index phase accuracy]
\label{sel:sharp-phase}
Uniformly for $\mathfrak a\in\mathfrak K$ and
$3\le n\le N_s^{\mathfrak a}$,
\begin{equation}
 \mathcal S_\lambda(\delta_n^{\mathfrak a})
 =\rho_c^{n-3}\{A_0L+\chi_\lambda(\mathfrak a)+r_{\mathfrak a}(\eps)\}
       +O(\eps^{3/32}),
 \qquad L=\log(1/\eps),
 \label{sel:sharp-phase-formula}
\end{equation}
where the entrance remainder satisfies
$\sup_{\mathfrak a\in\mathfrak K}|r_{\mathfrak a}(\eps)|\to0$.
The same comparison holds through any prescribed index $j=j(\eps)$
whose graph chain started at the actual base $\delta_3^{\mathfrak a}$ has
its $j$th base below $3s/2$, even if the observation time precedes $j$.
\end{lemma}
\begin{proof}
The uniform logarithmic entrance in \eqref{loc:log-entry} and
the exact-phase expansion give
\begin{equation}
 \mathcal S_\lambda(\delta_3^{\mathfrak a})
 =A_0L+\chi_\lambda(\mathfrak a)+r_{\mathfrak a}(\eps),
 \qquad
 r_{\mathfrak a}=r_{\mathfrak a}^{\log}+O(\eps^{3/32}),
 \label{sel:uniform-entry-phase}
\end{equation}
where the remainder tends to zero uniformly. We preserve its
parameter dependence by starting at the actual base.

Start the exact graph chain at the actual $\delta_3$, with fibre
$\Gamma_\lambda(\delta_3)$. The sharp third-entry estimate
\eqref{loc:third-entry} and the graph bound give
\begin{equation}
 E_3=B_3\le
 |\xi_3^r-\xi_{*,\lambda}|+
 |\Gamma_\lambda(\delta_3)-\xi_{*,\lambda}|
 =O(\eps^{3/32}).
 \label{sel:sharp-start}
\end{equation}
The equality $E_3=B_3$ uses the identical starting base. The
two-chain budget \eqref{path:two-chain-budget} gives
\[
 E_n\le q^{n-3}B_3+
 C\sum_{i=3}^{n-1}q^{n-1-i}\omega_i,
 \qquad
 \sup_i\omega_i\le C(\eps^{37/64}+\eps^{7/32}).
\]
The geometric sum is bounded by $8$. Thus
\begin{equation}
 E_n=O(\eps^{3/32})
 \label{sel:sharp-shadow}
\end{equation}
through every comparison within those domains, without a factor counting the returns.
All constants are common on $\mathfrak K$, since the entrance,
return and graph bounds are common. In particular the bound persists
at fixed positive scales.

Let $\bar N_{2s}$ be the first graph index whose base is at least
$2s$. Compare until $M=\min(N_s,\bar N_{2s})$. Before $M$ both
inputs lie below $2s<\Delta$ and the inner cutoffs; the estimate
also applies to their first output $M$ in the outer domain.
If $\bar N_{2s}<N_s$, the graph base at $M$ is at least $2s$ while
the actual base is less than $s$. Their logarithmic distance is then
at least $\log2$, contradicting \eqref{sel:sharp-shadow}. Thus the
comparison covers all $n\le N_s$. No further iteration at an output
outside the input domain has been used.

For the last assertion, generate the original trajectory until
$N_{2s}$, as permitted by Lemma~\ref{path:stopped-cascade}.
The graph bases through $j$ are below $3s/2$, since $F(d)>d$.
If the actual chain reached $2s$ at or before $j$, its first such
output would have logarithmic distance at least $\log(4/3)$ from
the graph base. This again contradicts \eqref{sel:sharp-shadow}.
The actual events through $j$ therefore exist and admit the stated
comparison on the underlying trajectory.

The graph phase at $n$ is exactly
$\rho_c^{n-3}\mathcal S_\lambda(\delta_3)$.
The two-sided bound \eqref{gr:phase-bilip} transfers the logarithmic
base error \eqref{sel:sharp-shadow} to phase. Substitution of
\eqref{sel:uniform-entry-phase} proves \eqref{sel:sharp-phase-formula}.
\end{proof}
